\begin{proof}
\textbf{Theorem (Jensen’s Inequality).}
Let $f:\mathbb R\to\mathbb R$ be convex.  
For real numbers $x_{1},\dots ,x_{n}$ and weights $\lambda_{1},\dots ,\lambda_{n}$ satisfying  
\[
\lambda_{i}\ge 0\;(i=1,\dots ,n),\qquad 
\sum_{i=1}^{n}\lambda_{i}=1,
\]
we have
\[
f\!\Bigl(\sum_{i=1}^{n}\lambda_{i}x_{i}\Bigr)\;\le\;
\sum_{i=1}^{n}\lambda_{i}f(x_{i}).\tag{1}
\]

\medskip
\textbf{1. Base case $n=2.}
Write $\lambda_{1}=t$ and $\lambda_{2}=1-t$ with $t\in[0,1]$.  
Convexity of $f$ gives
\begin{align*}
f\bigl(t x_{1}+(1-t)x_{2}\bigr)\le t\,f(x_{1})+(1-t)\,f(x_{2}).\tag{2}
\end{align*}
Taking $t=\lambda_{1}$ yields \eqref{1} for $n=2$.

\medskip
\textbf{2. Induction hypothesis.}
Assume that \eqref{1} holds for every collection of $k$ points with any non‑negative
weights summing to $1$, for all $k\le n-1$ (with $n\ge3$).

\medskip
\textbf{3. Induction step (the case of $n$ points).}
Let $\lambda_{1},\dots ,\lambda_{n}$ be as in the statement and set
\[
\mu:=\sum_{i=1}^{n-1}\lambda_{i}.
\]

\medskip
\textbf{3.1 Degenerate sub‑case $\mu=0$.}
If $\mu=0$ then $\lambda_{i}=0$ for $i=1,\dots ,n-1$ and $\lambda_{n}=1$, so both
sides of \eqref{1} reduce to $f(x_{n})$ and the inequality holds with equality.

\medskip
\textbf{3.2 Non‑degenerate sub‑case $\mu>0$.}
Define normalised weights for the first $n-1$ points
\[
\alpha_{i}:=\frac{\lambda_{i}}{\mu},\qquad i=1,\dots ,n-1,
\]
so that $\alpha_{i}\ge0$ and $\sum_{i=1}^{n-1}\alpha_{i}=1$.
Set
\[
y:=\sum_{i=1}^{n-1}\alpha_{i}x_{i}.
\]
Then $\mu y=\sum_{i=1}^{n-1}\lambda_{i}x_{i}$.

Applying the induction hypothesis to the points $(x_{i})_{i=1}^{n-1}$ with
weights $(\alpha_{i})$ yields
\begin{align*}
f(y)\le \sum_{i=1}^{n-1}\alpha_{i}f(x_{i}).\tag{3}
\end{align*}
Multiplying \eqref{3} by the positive scalar $\mu$ gives
\begin{align*}
\mu f(y)\le \sum_{i=1}^{n-1}\lambda_{i}f(x_{i}).\tag{4}
\end{align*}

Since $\mu+\lambda_{n}=1$ and $\mu\in[0,1]$, convexity applied to the two points
$y$ and $x_{n}$ with weights $\mu$ and $\lambda_{n}$ yields
\begin{align*}
f\!\bigl(\mu y+\lambda_{n}x_{n}\bigr)\le
\mu f(y)+\lambda_{n}f(x_{n}).\tag{5}
\end{align*}
Because $\mu y=\sum_{i=1}^{n-1}\lambda_{i}x_{i}$, the left–hand side of \eqref{5}
is $f\!\bigl(\sum_{i=1}^{n}\lambda_{i}x_{i}\bigr)$.  Substituting \eqref{4} into the
right–hand side of \eqref{5} we obtain
\[
f\!\Bigl(\sum_{i=1}^{n}\lambda_{i}x_{i}\Bigr)
\le
\sum_{i=1}^{n-1}\lambda_{i}f(x_{i})+\lambda_{n}f(x_{n})
=
\sum_{i=1}^{n}\lambda_{i}f(x_{i}),
\]
which is exactly \eqref{1} for $n$ points.  This completes the induction.

\medskip
\textbf{4. Equality condition.}
Equality in \eqref{1} can occur only when equality holds in every convexity
application used above.

\begin{itemize}
\item In the base case $n=2$, equality in \eqref{2} holds precisely when either
  $t=0$, $t=1$ (i.e. one weight is zero) \emph{or} $f$ is affine on the segment
  $[x_{1},x_{2}]$ (the latter includes $x_{1}=x_{2}$ when $f$ is strictly convex).

\item In the induction step, the second convexity application \eqref{5} is an
  equality iff either $\mu=0$ or $\lambda_{n}=0$ (so the combination reduces to a
  single point) \emph{or} $f$ is affine on the segment $[y,x_{n}]$.
\end{itemize}
Consequently, equality in \eqref{1} holds exactly in one of the following
situations:

\begin{enumerate}
\item \textbf{Trivial concentration of weight:} there exists $k$ with
  $\lambda_{k}=1$ and $\lambda_{i}=0$ for $i\neq k$; both sides equal $f(x_{k})$.

\item \textbf{All positively‑weighted points coincide:} if $\lambda_{i}>0$ and
  $\lambda_{j}>0$ then $x_{i}=x_{j}$.  The convex combination reduces to that common
  value, giving equality.

\item \textbf{Affineness of $f$ on the relevant convex hull:} the restriction of
  $f$ to the convex hull of $\{x_{i}:\lambda_{i}>0\}$ is affine (linear).  This
  includes the case where $f$ is globally affine.
\end{enumerate}
These possibilities are exhaustive; each indeed yields equality in \eqref{1}.

\medskip
Thus for any convex $f$, any real numbers $x_{1},\dots ,x_{n}$, and any
non‑negative weights $\lambda_{1},\dots ,\lambda_{n}$ summing to $1$, Jensen’s
inequality \eqref{1} holds, with equality precisely in the cases described
above. \qed
\end{proof}